\chapter{Inequalities and Absolute Values}

Solve each. Mind the sign when multiplying by a negative or by an
expression of unknown sign.

\section{Core Rules}

\begin{rulebox}{The one rule that separates inequalities from equations}
Adding or subtracting anything is safe. Multiplying or dividing by a
\textbf{positive} number is safe. Multiplying or dividing by a
\textbf{negative} number \textbf{reverses} the inequality:
\[
  -2x<6\quad\Longrightarrow\quad x>-3 .
\]
Consequently you may never multiply by an expression whose sign is unknown.
$x+1$, $x-2$, $x$ itself: all change sign. Multiplying an inequality by
$x+1$ is not a step, it is two steps with two different answers.
\end{rulebox}

\begin{rulebox}{Rational and quadratic inequalities: the sign table}
Bring everything to one side, so the question becomes ``where is this
expression $\le0$?'', and factor:
\[
  \frac{P(x)}{Q(x)}\le0 .
\]
Mark on the number line every zero of $P$ (candidate solutions, included for
$\le$) and every zero of $Q$ (poles, \textbf{always excluded}). These points
cut the line into intervals, and the expression keeps one sign on each
interval, so testing a single convenient point per interval settles it.
\end{rulebox}

\begin{rulebox}{Absolute value: distance, and the two standard shapes}
$|x|$ is the distance from $x$ to $0$; $|x-a|$ is the distance from $x$
to~$a$. Formally $|x|=x$ for $x\ge0$ and $|x|=-x$ for $x<0$. For $a>0$:
\[
  |X|<a \iff -a<X<a
  \qquad\text{(one interval, an ``and'')},
\]
\[
  |X|>a \iff X<-a\ \text{ or }\ X>a
  \qquad\text{(two rays, an ``or'')}.
\]
$|X|=a$ splits into $X=a$ or $X=-a$; $|X|=0$ only at $X=0$; $|X|<0$ never.
\end{rulebox}

\begin{rulebox}{Two absolute values, and the triangle inequality}
$|A|\le|B|$ with both sides non-negative may be \textbf{squared}:
$A^{2}\le B^{2}$, which is a quadratic inequality. Squaring is legitimate
here only because both sides are known to be $\ge0$. Finally, for all real
$a,b$,
\[
  |a+b|\le|a|+|b| ,
\]
with equality exactly when $a$ and $b$ have the same sign (or one is zero).
The reverse inequality is false, and so is equality in general.
\end{rulebox}

\section{Worked Examples}

\begin{worked}{(a) Solve $\dfrac{3}{x+1}\le 2$}
\wstep{Step 1. Do \emph{not} multiply by $x+1$.}
The obvious move gives $3\le2(x+1)$, hence $x\ge\tfrac12$. It is wrong,
because $x+1$ is negative for $x<-1$ and multiplying by it would flip the
sign. Test $x=-2$ in the original: $3/(-1)=-3\le2$. \checkmark\ So $x=-2$
solves the inequality but is missing from that answer.

\wstep{Step 2. Move everything to one side instead.}
\[
  \frac{3}{x+1}-2\le0
  \;\Longrightarrow\;
  \frac{3-2(x+1)}{x+1}\le0
  \;\Longrightarrow\;
  \frac{1-2x}{x+1}\le0 .
\]
No sign assumption was made, because nothing was multiplied.

\wstep{Step 3. Locate the critical points.}
The numerator vanishes at $x=\tfrac12$; the denominator vanishes at $x=-1$.
These split the line into $(-\infty,-1)$, $(-1,\tfrac12)$,
$(\tfrac12,\infty)$.

\wstep{Step 4. Test one point in each interval.}
\begin{center}\footnotesize
\begin{tabular}{lccc}
  \hline
  interval & test $x$ & $\dfrac{1-2x}{x+1}$ & sign\\
  \hline
  $(-\infty,-1)$ & $-2$ & $5/(-1)$ & $-$ \\
  $(-1,\tfrac12)$ & $0$ & $1/1$ & $+$ \\
  $(\tfrac12,\infty)$ & $1$ & $-1/2$ & $-$ \\
  \hline
\end{tabular}
\end{center}

\wstep{Step 5. Collect, and decide the endpoints separately.}
We need ``$\le0$'', so the two negative intervals qualify. At $x=\tfrac12$
the expression is $0$, and $0\le0$, so $\tfrac12$ is \emph{included}. At
$x=-1$ the expression is undefined, so $-1$ is \emph{excluded}:
\[
  \boxed{x<-1\quad\text{or}\quad x\ge\tfrac12}
\]
\wtrap{Zeros of the numerator are included when the inequality is weak;
zeros of the denominator are excluded always. Treating both alike is the
standard slip, and it is the difference between a bracket and a parenthesis
in the answer.}
\end{worked}

\begin{worked}{(b) Solve $|2x-1|\le|x+4|$}
\wstep{Step 1. Both sides are non-negative, so squaring is safe.}
An absolute value is never negative, so $|A|\le|B|$ and $A^{2}\le B^{2}$ say
the same thing. This avoids the four-case analysis entirely.
\[
  (2x-1)^{2}\le(x+4)^{2} .
\]
\wtrap{Squaring is only legitimate because \emph{both} sides are known
non-negative. Squaring $2x-1\le x+4$ would be invalid --- the left side can
be negative, and squaring would reverse the comparison.}

\wstep{Step 2. Expand and collect on one side.}
\[
  4x^{2}-4x+1\le x^{2}+8x+16
  \;\Longrightarrow\;
  3x^{2}-12x-15\le0
  \;\Longrightarrow\;
  x^{2}-4x-5\le0 ,
\]
dividing by the positive number $3$, which does not flip anything.

\wstep{Step 3. Factor and read the sign.}
$x^{2}-4x-5=(x-5)(x+1)$. This upward parabola is $\le0$ exactly
\emph{between} its roots:
\[
  \boxed{-1\le x\le 5}
\]
\wstep{Step 4. Check the ends and one interior point.}
$x=-1$: $|-3|=3$, $|3|=3$, equal. \checkmark\quad
$x=5$: $|9|=9$, $|9|=9$, equal. \checkmark\quad
$x=0$: $|-1|=1\le|4|=4$. \checkmark\quad
$x=6$: $|11|=11\le|10|=10$ is false, correctly outside.
\wtrap{An upward parabola is negative \emph{between} its roots and positive
outside; a downward one is the reverse. Reading this off the wrong way round
turns the answer inside out, and the check in Step~4 catches it in seconds.}
\end{worked}

\section{Practice Problems}

\begin{problems}
  \item $7-3x\le 1$
  \item $\displaystyle\frac{2x+1}{3}>5$
  \item $5-2(x-1)\ge 3$
  \item $-1\le 4-x\le 3$
  \item $\displaystyle\frac{x-1}{2}<\frac{x+1}{3}$
  \item $0<2x+6<10$
  \item Phillips curve: $\pi_{t}=2+0.5\,u_{t}$. Solve $\pi_{t}\le 3$ for
        $u_{t}$.
  \item LP feasibility: $x+y\le 10,\ x\ge 0,\ y\ge 0$. Express the range of
        $y$ for a fixed $x\in[0,10]$.
  \item $(x-2)(x+3)>0$
  \item $x^{2}-x-12\le 0$
  \item Show that $x^{2}+4>0$ for all $x\in\R$.
  \item $-2x^{2}+8\ge 0$
  \item $\displaystyle\frac{1}{x}>0$
  \item $\displaystyle\frac{x+1}{x-2}>0$
  \item $|3x|=12$
  \item $|x+7|=0$
  \item $|2x+5|-3=0$
  \item For which $x$ is $|x|=x$?
  \item $|x-5|=|x+3|$ (use the geometric interpretation)
  \item $|x|+|x-2|=2$
  \item $|x-1|>|x+1|$
  \item $|x|<2|x-3|$
  \item $|2x-6|=4$
  \item $|2x-1|\le|x+4|$
  \item Which is true for all $a,b\in\R$:
        (a) $|a+b|=|a|+|b|$;
        (b) $|a+b|\le|a|+|b|$;
        (c) $|a+b|\ge|a|+|b|$.
  \item $|x-3|>0$
  \item $|x|^{2}-4|x|+3=0$ (Hint: let $u=|x|$.)
  \item Budget: $p_{x}x+p_{y}y\le I,\ x=0$. State the max $y$.
  \item $\displaystyle\frac{3}{x+1}\le 2$ (beware sign of $x+1$)
  \item $x^{2}\ge 2x$
\end{problems}

\section{Solutions}

\begin{solutions}
  \item $x\ge 2$
  \item $x>7$
  \item $x\le 2$
  \item $1\le x\le 5$
  \item $x<5$
  \item $-3<x<2$
  \item $u_{t}\le 2$
  \item $0\le y\le 10-x$
  \item $x>2$ or $x<-3$
  \item $-3\le x\le 4$
  \item Since $x^{2}\ge 0$, we have $x^{2}+4\ge 4>0$.
  \item $-2\le x\le 2$
  \item $x>0$
  \item $x>2$ or $x<-1$
  \item $x=\pm 4$
  \item $x=-7$
  \item $x=-1,\,-4$
  \item $x\ge 0$
  \item $x=1$ (midpoint of $5$ and $-3$)
  \item $0\le x\le 2$ (any point between $0$ and $2$ has distances summing
        to $2$)
  \item $x<0$
  \item $x<2$ or $x>6$
  \item $x=5,\,1$
  \item $-1\le x\le 5$
  \item (b): the triangle inequality; (a) holds only when $a,b$ share a
        sign.
  \item $x\neq 3$
  \item $|x|=1$ or $|x|=3\Rightarrow x=\pm1,\pm3$
  \item $y\le I/p_{y}$
  \item $x<-1$ or $x\ge\tfrac12$
  \item $x\le 0$ or $x\ge 2$
\end{solutions}
